\section{Evaluación, gobernanza y despliegue}

\subsection{Matriz de evaluación multinivel}

Un proceso de evaluación completo debe incluir benchmark automático, trajectory review y stress test:

\begin{itemize}
    \item Métricas numéricas como `score`, `TD error`, `episode length`
    \item Revisión manual de failure trajectory / reward hack
    \item Inyección de ruido, belief shift y perturbation del entorno para probar la estabilidad
\end{itemize}

\begin{table}[H]
\centering
\begin{tabular}{p{0.25\textwidth} p{0.25\textwidth} p{0.2\textwidth}}
\toprule
Nivel & Métrica & Estrategia de activación \\
\midrule
Training eval & Q target, loss, entropy & Divergence or reward plateau \\
Safety audit & reward hack detection, adversarial action & anomalous trajectories \\
Deployment sim & latency, repeatability, replicates & hardware change / rollout \\
\bottomrule
\end{tabular}
\caption{Matriz de evaluación y gobernanza}
\end{table}

\begin{knowledgebox}{Recomendaciones de monitoreo de logs}
Registrar TD error distribution, max Q target, exploration ratio; configurar \texttt{td\_error > threshold} para activar replay augmentation, y \texttt{reward variance > 2σ} para activar reward shaping review.
\end{knowledgebox}

\subsection{Checklist de despliegue y gobernanza}

\begin{warningbox}{Las tres «puertas» antes de salir a producción}
\begin{enumerate}
    \item El replay buffer solo puede contener transitions verificadas mediante human-in-the-loop.
    \item El checkpoint debe evitar reward hack, por ejemplo mediante reward clipping + sanity check trajectories.
    \item Definir condiciones claras de rollback (loss divergence, reward drop > threshold).
\end{enumerate}
\end{warningbox}

\begin{figure}[H]
\centering
\includegraphics[width=0.86\textwidth,page=2]{06_cs224r_qlearning_2025.pdf}
\caption{Slide: Q-Learning deployment guardrails y monitoring stack}
\end{figure}

\subsection{Evidence Matrix y línea de tiempo de acciones}
\subsection{Evidence Matrix y línea de tiempo de acciones}

Registrar el evidence matrix en cada ronda de entrenamiento ayuda a ubicar rápidamente la failure root cause. La misma matriz también puede mapearse a un action timeline:

\begin{table}[H]
\centering
\begin{tabular}{p{0.28\textwidth} p{0.58\textwidth}}
\toprule
Evidence & Action \\
\midrule
Prediction explosion & Retroceder a un checkpoint con learning rate bajo + aumentar el target damping \\
Reward spike & Revisar los reward shaping logs, pausar el training y revisar el newly added shaping term \\
TD error drift & Activar replay augmentation y aumentar el PER alpha \\
Policy collapse (success rate drop) & Enforce early rollback + high-priority fail buffer replay \\
\bottomrule
\end{tabular}
\caption{Evidence matrix y su línea de tiempo de acciones correspondiente}
\end{table}

\begin{knowledgebox}{Actions should be traceable}
Registrar el context (checkpoint id, evidence source, actor responsible) de cada triggered action facilita la auditoría y el replay review.
\end{knowledgebox}

\subsection{Resumen de la sección}

La matriz de evaluación y el checklist de gobernanza garantizan que el algoritmo de Q-Learning no se vea afectado por reward hack o value explosion después del despliegue.
